% generated: gpt-6 2026-10-29 role=assisted; see PROVENANCE.md
\documentclass[11pt]{article}
\usepackage[margin=25mm]{geometry}
\usepackage{latex/clean}
\setlength{\emergencystretch}{3em}
\title{Magnetic matrix statements in Clean}
\author{}
\date{29 October 2026}
\begin{document}
\maketitle
\noindent{\footnotesize LLM-assisted specimen. Source attributions retain their
original authorship; see PROVENANCE.md.}\par\medskip
These six statements have Lean proofs in \path{studies/magnetic/Verified.lean}.
The Laplacian uses the unnormalized convention $L_q=D-M_q$, with
$D_{ii}=\sum_j(A_{ij}+A_{ji})$. The inverse endpoint result is a counterexample
to the displayed odd-length specialization of v1 equation (24).
The matrix $E_{01}$ is $\left(\begin{smallmatrix}0&1\\0&0\end{smallmatrix}\right)$.
% generated: gpt-6 2026-10-29 role=assisted; LLM-authored portions; see PROVENANCE.md
% Generated statements, not proof certificates.
\begin{cleanresult}{magnetic powers encode walk profiles}\cleanattribution{\href{https://arxiv.org/html/2506.07343v1}{[Huang, Gleich, Li 2025, equation (1), Supplement S.II equation (17), Theorem 1]}}
Let $n$ be a natural number.
Let $A$ be a matrix of size $n\times n$ over $\mathbb{R}$.
Let $m$ be a natural number.
Let $q$ be a real number.
Then \[{M_{q}(A)}^{m} = e^{2\pi i -(q \, m)} \, \operatorname{sumBefore}(m + 1,\left(k \mapsto e^{2\pi i 2 \, q \, k} \, \Phi_{A}(m,k)\right)).\]
\par\smallskip{\footnotesize Adapted from the cited source. Scope: forward transform only.}
\end{cleanresult}

\begin{cleanresult}{magnetic adjacency is Hermitian}\cleanattribution{\href{https://arxiv.org/html/2506.07343v1}{[Huang, Gleich, Li 2025, magnetic graph matrix definition]}}
Let $n$ be a natural number.
Let $A$ be a matrix of size $n\times n$ over $\mathbb{R}$.
Let $q$ be a real number.
Then \[M_{q}(A)\text{ is Hermitian}.\]
\par\smallskip{\footnotesize Adapted from the cited source.}
\end{cleanresult}

\begin{cleanresult}{magnetic Laplacian is Hermitian}\cleanattribution{[Clean magnetic formalization, unnormalized D minus M convention]}
Let $n$ be a natural number.
Let $A$ be a matrix of size $n\times n$ over $\mathbb{R}$.
Let $q$ be a real number.
Then \[L_{q}(A)\text{ is Hermitian}.\]
\end{cleanresult}

\begin{cleanresult}{zero potential symmetrizes adjacency}\cleanattribution{\href{https://arxiv.org/html/2506.07343v1}{[Huang, Gleich, Li 2025, discussion following equation (1)]}}
Let $n$ be a natural number.
Let $A$ be a matrix of size $n\times n$ over $\mathbb{R}$.
Then \[M_{0}(A) = A + A^{\mathsf{T}}.\]
\par\smallskip{\footnotesize Adapted from the cited source. Scope: first power.}
\end{cleanresult}

\begin{cleanresult}{two step mixed walk profile}\cleanattribution{\href{https://arxiv.org/html/2506.07343v1}{[Huang, Gleich, Li 2025, walk profile examples preceding equation (1)]}}
Let $n$ be a natural number.
Let $A$ be a matrix of size $n\times n$ over $\mathbb{R}$.
Then \[\Phi_{A}(2,1) = A \, A^{\mathsf{T}} + A^{\mathsf{T}} \, A.\]
\par\smallskip{\footnotesize Adapted from the cited source.}
\end{cleanresult}

\begin{cleanresult}{odd length inverse endpoint counterexample}\cleanattribution{[Clean magnetic formalization, inverse\_v1\_length\_one\_fails]}
Then \[\Phi_{E_{01}}(1,1) \ne \frac{1}{2} \, M_{0}(E_{01}).\]
\par\smallskip{\footnotesize Scope: v1 equation (24) at m = 1.}
\end{cleanresult}

\section*{References}
\begin{description}
\item[Huang, Gleich, Li 2025] Yinan Huang, David F. Gleich, Pan Li. Powers of Magnetic Graph Matrix: Fourier Spectrum, Walk Compression, and Applications. Version arXiv:2506.07343v1. \href{https://arxiv.org/html/2506.07343v1}{Source}
\item[Clean magnetic formalization] . Supporting results for the magnetic graph matrix.
\end{description}

\end{document}

% generated: gpt-6 2026-10-29 role=assisted; LLM-authored portions; see PROVENANCE.md
